\cvsection{Expériences professionnelles}

\begin{cventries}
  \cventry
    {Stagiaire ingénieur de recherche en IA}
    {INRIA}
    {Bordeaux, France}
    {Mars 2026 -- Septembre 2026}
    {
      \begin{cvitems}
        \item {Analyse du besoin des chercheurs et développement en autonomie d’un assistant IA pour la recherche d’articles scientifiques.}
        \item {Conception de prompts pour adapter le comportement de l’agent à la recherche documentaire et guider l’utilisation de ses outils.}
        \item {Développement d’un agent multi-étapes sans framework d’orchestration, avec intégration de l’API OpenAI et gestion des appels d’outils, de l’état et de la mémoire.}
        \item {Développement d’API Python/FastAPI et d’une interface React/Next.js permettant à l’utilisateur de choisir les messages conservés dans le contexte du modèle.}
        \item {Mise en œuvre d’un pipeline RAG avec chunking, embeddings et recherche sémantique dans ChromaDB.}
        \item {Définition d’un protocole de comparaison de huit LLM sur 70 articles, avec mesure de la réussite de la recherche selon la précision des requêtes et comparaison de la latence d’inférence.}
        \item {Évaluation d’OpenClaw dans un environnement Docker isolé pour étudier les risques liés à l’accès d’un agent aux outils et aux fichiers.}
        \item {Déploiement de LLM via vLLM sur cluster HPC équipé de GPU NVIDIA et automatisation de l’intégration avec Docker et GitHub CI/CD.}
      \end{cvitems}
    }

  \cventry
    {Stagiaire ingénieur de recherche en apprentissage par renforcement}
    {Université Karunya}
    {Coimbatore, Inde}
    {Juin 2025 -- Septembre 2025}
    {
      \begin{cvitems}
        \item {Développement d’un modèle d’apprentissage par renforcement pour la navigation autonome d’un robot hospitalier, puis déploiement et tests sur Raspberry Pi.}
      \end{cvitems}
    }
\end{cventries}
